\documentclass[11pt]{article}

\usepackage[margin=1in]{geometry}
\usepackage{amsmath,amssymb}
\usepackage[hidelinks]{hyperref}

\title{Status note for arXiv:1005.3634, Problem 1}
\author{fa\_banach\_001, agent\_lane\_15}
\date{2026-06-28}

\begin{document}
\maketitle

\section*{Status}

This is a lightweight literature-status packet.  The extracted problem is
\emph{already answered in later literature}; no new proof is claimed here.

\section*{Source question}

The source paper is F. Bermudez, A. Bonilla, F. Martinez-Gimenez, and A. Peris,
\emph{Li-Yorke and distributionally chaotic operators}, J. Math. Anal. Appl.
373 (2011), 83--93, DOI
\href{https://doi.org/10.1016/j.jmaa.2010.06.011}{10.1016/j.jmaa.2010.06.011},
arXiv:1005.3634.

In Section 4, ``Distributionally chaotic operators'', immediately after
Proposition 16, the paper asks:
\[
\text{Does every distributionally chaotic operator } T:X\to X
\text{ admit a distributionally irregular vector?}
\]
In the arXiv source this is Problem 1.

\section*{Answering theorem}

The answer source is N. C. Bernardes Jr., A. Bonilla, V. Muller, and A. Peris,
\emph{Distributional chaos for linear operators}, J. Funct. Anal. 265 (9)
(2013), 2143--2163, DOI
\href{https://doi.org/10.1016/j.jfa.2013.06.019}{10.1016/j.jfa.2013.06.019}.

The relevant result is Theorem 12 of the 2013 JFA paper.  Later accessible
sources restate it in the following form: for a bounded linear operator
\(T\) on a Banach space \(X\), the following conditions are equivalent:
\(T\) satisfies the Distributional Chaos Criterion, \(T\) has a
distributionally irregular vector, \(T\) is distributionally chaotic, and
\(T\) admits a distributionally chaotic pair.  Thus the implication
\[
T \text{ distributionally chaotic}
\quad\Longrightarrow\quad
T \text{ has a distributionally irregular vector}
\]
is exactly the affirmative answer to Problem 1.

\section*{Provenance check}

Crossref metadata for the 2013 JFA article records a reference to the 2011
source article by DOI
\href{https://doi.org/10.1016/j.jmaa.2010.06.011}{10.1016/j.jmaa.2010.06.011},
title, journal, volume, issue, first page, and year.  This supports placement
in \texttt{literature\_already\_answered/}: the later answer source is
separate from, and bibliographically connected to, the source problem paper.

The packet includes the source arXiv PDF and Crossref metadata for both DOIs.
The decisive 2013 JFA paper does not appear as a local arXiv PDF in the corpus;
therefore the packet also includes the accessible witness arXiv:1807.05191,
which restates Bernardes--Bonilla--Muller--Peris Theorem 12 as Theorem 1.1.

\section*{Conclusion}

The target arXiv:1005.3634 is not viable for a new full proof or
counterexample on this question.  Problem 1 is already solved in the
affirmative by Bernardes--Bonilla--Muller--Peris, Theorem 12.

\begin{thebibliography}{9}

\bibitem{BBMP2011}
F. Bermudez, A. Bonilla, F. Martinez-Gimenez, and A. Peris,
\emph{Li-Yorke and distributionally chaotic operators},
J. Math. Anal. Appl. 373 (2011), 83--93.

\bibitem{BBMP2013}
N. C. Bernardes Jr., A. Bonilla, V. Muller, and A. Peris,
\emph{Distributional chaos for linear operators},
J. Funct. Anal. 265 (9) (2013), 2143--2163.

\bibitem{Witness2018}
C. Chen,
\emph{Distributional chaos for weighted translation operators on groups},
arXiv:1807.05191.

\end{thebibliography}

\end{document}
